\documentclass[11pt,a4paper]{article}

\usepackage[margin=1in]{geometry}
\usepackage[T1]{fontenc}
\usepackage[utf8]{inputenc}
\usepackage{lmodern}
\usepackage{setspace}
\usepackage{booktabs}
\usepackage{longtable}
\usepackage{array}
\usepackage{enumitem}
\usepackage{hyperref}
\usepackage{xcolor}
\usepackage{graphicx}
\usepackage{float}

\hypersetup{
  colorlinks=true,
  linkcolor=blue,
  urlcolor=blue,
  citecolor=blue
}

\setlength{\parskip}{0.6em}
\setlength{\parindent}{0pt}
\onehalfspacing

\title{\textbf{Comprehensive Technical Report}\\
\large Deciphering the Molecular Grammar of Hypertrophic Cardiomyopathy (HCM)}
\author{HCM Project Team \\ RVCE Experiential Learning (2025--26)}
\date{\today}

\begin{document}

\maketitle
\thispagestyle{empty}

\begin{abstract}
This report documents the complete technical lifecycle of the HCM pathogenicity prediction project and expands the repository narrative into an end-to-end engineering and research record. We cover project motivation, data governance, leakage prevention strategy, feature and embedding design, model architectures, cross-gene validation, calibration analysis, in silico mutagenesis outputs, external benchmarking against EVE, VUS re-stratification outcomes, and reproducibility operations. The repository delivers a practical and auditable pipeline that combines tabular structural/biophysical features with ESM-2 embedding deltas in a Two-Tower neural framework, and compares its performance to a robust Random Forest baseline.
\end{abstract}

\newpage
\tableofcontents
\newpage

\section{Executive Summary}
Hypertrophic Cardiomyopathy (HCM) has a large unresolved interpretation burden due to variants categorized as \textit{Variant of Uncertain Significance} (VUS). This project addresses the problem by building a zero-leakage machine learning workflow focused on pathogenicity prediction in nine sarcomeric genes.

\textbf{Core outcomes from repository artifacts:}
\begin{itemize}[leftmargin=1.2cm]
    \item A complete three-phase implementation pipeline (Month 1, Month 2, Month 3) is available in \texttt{scripts/}.
    \item Labeled training cohort contains \textbf{1,954} variants with a pathogenic-majority class distribution.
    \item Full variant cohort contains \textbf{8,137} variants, including \textbf{4,523} VUS used for re-stratification.
    \item Two production model families are shipped: \textbf{Two-Tower Hybrid} and \textbf{Baseline Random Forest}.
    \item Per-gene LOGO metrics, calibration plots, and ISM heatmaps are persisted for audit and downstream interpretation.
    \item Ranked VUS pathogenicity lists are generated and exported for translational review workflows.
\end{itemize}

\section{Project Context and Clinical Motivation}
\subsection{Problem Setting}
HCM-associated missense variants are abundant in clinical sequencing datasets, but many observed substitutions do not have sufficient functional or population evidence to support unambiguous benign/pathogenic classification. This creates clinical uncertainty, delays actionability, and increases dependence on manual curation.

\subsection{Why a Computational Prioritization Layer Is Needed}
A model with calibrated probabilities does not replace expert adjudication; instead, it provides an ordered decision-support signal to prioritize expensive downstream validation. In this project, the computational layer is optimized for:
\begin{itemize}[leftmargin=1.2cm]
    \item \textbf{Generalization across genes} rather than memorization of gene-specific priors,
    \item \textbf{Leakage resistance} to avoid inflation from known clinical annotations,
    \item \textbf{Reusability} via persisted embeddings and model artifacts,
    \item \textbf{Interpretability support} through mutation landscape visual outputs.
\end{itemize}

\subsection{Project Objectives}
The implementation objectives are:
\begin{enumerate}[leftmargin=1.2cm]
    \item Train robust pathogenicity classifiers using both engineered and representation-learned signals.
    \item Evaluate by \textit{Leave-One-Gene-Out} (LOGO) to assess cross-gene transfer performance.
    \item Apply probability calibration for better score reliability.
    \item Generate ranked VUS outputs from full-variant cohorts.
    \item Produce reproducible artifacts and visualizations suitable for technical reporting.
\end{enumerate}

\section{Repository Scope and Deliverables}
\subsection{Top-Level Structure}
\begin{longtable}{p{0.30\textwidth} p{0.65\textwidth}}
\toprule
\textbf{Path} & \textbf{Role in pipeline} \\
\midrule
\texttt{data/HCM\_labeled\_final.csv} & Curated labeled cohort for model training and LOGO evaluation. \\
\texttt{data/HCM\_all\_variants\_v2.csv} & Extended cohort used to isolate VUS and run re-stratification. \\
\texttt{data/esm2\_delta\_embeddings.npy} & Persisted embedding deltas generated in Month 1 for downstream reuse. \\
\texttt{scripts/execute\_month1.py} & Data cleaning, leaky field removal, ESM extraction, baseline LOGO RF step. \\
\texttt{scripts/execute\_month2.py} & Two-Tower training, LOGO loop, calibration, bootstrap CIs, ISM generation. \\
\texttt{scripts/execute\_month3.py} & Final model training on full labeled data and VUS ranking generation. \\
\texttt{models/} & Final artifacts: Two-Tower weights, RF model bundle, Platt calibrator pickle. \\
\texttt{results/} & Model-level and VUS-level output CSVs for downstream analysis. \\
\texttt{figures/} & Calibration and in silico mutagenesis visual outputs for each target gene. \\
\texttt{eve\_benchmark/} & External benchmark scripts and outputs against EVE reference scores. \\
\texttt{paper/} & Manuscripts and this expanded LaTeX report source. \\
\bottomrule
\end{longtable}

\subsection{Operational Deliverables Produced by the Pipeline}
\begin{itemize}[leftmargin=1.2cm]
    \item Trained Two-Tower neural weights: \texttt{models/hcm\_final\_two\_tower\_model.pth}
    \item Random Forest model bundle: \texttt{models/hcm\_final\_rf\_model.joblib}
    \item Calibration model: \texttt{models/hcm\_platt\_calibrator.pkl}
    \item LOGO metrics table: \texttt{results/logo\_metrics.csv}
    \item VUS ranking outputs: \texttt{results/VUS\_restratification\_two\_tower.csv}, \texttt{results/VUS\_restratification\_rf.csv}
    \item Calibration plots and per-gene ISM heatmaps under \texttt{figures/}
\end{itemize}

\section{Data Assets and Cohort Summary}
\subsection{Labeled Cohort Characteristics}
\begin{itemize}[leftmargin=1.2cm]
    \item Total labeled variants: \textbf{1,954}
    \item Pathogenic class: \textbf{1,441} (73.7\%)
    \item Benign class: \textbf{513} (26.3\%)
    \item Genes covered: \textbf{9} (\texttt{MYH7, MYBPC3, TNNT2, TNNI3, TPM1, ACTC1, MYL2, MYL3, TNNC1})
\end{itemize}

\subsection{Per-Gene Variant Counts (Labeled)}
\begin{center}
\begin{tabular}{lr}
\toprule
\textbf{Gene} & \textbf{Count} \\
\midrule
ACTC1 & 64 \\
MYBPC3 & 429 \\
MYH7 & 986 \\
MYL2 & 54 \\
MYL3 & 37 \\
TNNC1 & 36 \\
TNNI3 & 118 \\
TNNT2 & 130 \\
TPM1 & 100 \\
\bottomrule
\end{tabular}
\end{center}

\subsection{Full Cohort for Re-stratification}
\begin{itemize}[leftmargin=1.2cm]
    \item Total variants in full set: \textbf{8,137}
    \item Variants marked uncertain significance: \textbf{4,523}
\end{itemize}

\section{Leakage Prevention and Data Governance}
\subsection{Leakage Risk in Clinical Variant Learning}
Clinical variant datasets often contain meta-annotations that encode outcome-related curation information. If such attributes are fed to the model, estimated performance can become artificially inflated and not reflective of biological generalization.

\subsection{Implemented Leakage Guard}
The implementation explicitly removes the following columns when present:
\begin{center}
\begin{tabular}{ll}
\toprule
\textbf{Removed field} & \textbf{Reason for exclusion} \\
\midrule
\texttt{pop\_freq} & May encode prior pathogenicity assumptions through rarity patterns \\
\texttt{disease} & Directly tied to diagnostic interpretation context \\
\texttt{sources} & Curation provenance may correlate with label reliability \\
\texttt{genomic\_loc} & Non-generalizable metadata with leakage potential \\
\texttt{review\_status} & Human review confidence can proxy target labels \\
\texttt{clin\_sig} & Directly linked to target interpretation semantics \\
\bottomrule
\end{tabular}
\end{center}

\subsection{Engineering Implication}
This explicit drop-step appears in all three execution phases and is central to the ``zero-leakage'' claim. It ensures that both intermediate validation and final VUS scoring use feature spaces that are less contaminated by label-adjacent metadata.

\section{Feature Engineering and Representation Design}
\subsection{Tabular/Biophysical Feature Block}
The tabular tower receives numeric/boolean features that include structural context flags, physicochemical descriptors, and mutation-derived encodings. From script behavior and excluded fields, this block is assembled by filtering non-identifiers and retaining numeric/boolean columns after leakage removal.

\subsection{Sequence Representation Block}
The sequence tower consumes ESM-derived embedding deltas:\\
\textit{mutant representation} $-$ \textit{wild-type representation}, aggregated from token-level outputs by mean pooling.

\subsection{Mutation-aware Sequence Construction}
For 11-mer windows, mutation replacement is centered at index 5 when the reference residue matches the expected center residue. A fallback string replacement is used otherwise, enabling robust handling when windows are irregular.

\subsection{Combined Signal Hypothesis}
The modeling assumption is that:
\begin{itemize}[leftmargin=1.2cm]
    \item tabular descriptors capture explicit structural/biophysical priors,
    \item ESM deltas capture context-sensitive sequence perturbation effects,
    \item fusion of both yields better transferability than either source alone.
\end{itemize}

\section{Modeling Pipeline by Development Phase}
\subsection{Month 1: Foundation Layer}
\textbf{Implemented tasks:}
\begin{itemize}[leftmargin=1.2cm]
    \item Leakage-aware cleaning of labeled data,
    \item ESM-2 delta embedding extraction and persistence,
    \item Baseline Random Forest training under a LOGO hold-out setting,
    \item Serialization of baseline model artifact.
\end{itemize}

\textbf{Primary outputs:}
\begin{itemize}[leftmargin=1.2cm]
    \item \texttt{data/esm2\_delta\_embeddings.npy}
    \item \texttt{models/hcm\_logo\_baseline\_model.joblib}
\end{itemize}

\subsection{Month 2: Hybrid Modeling and Evaluation Layer}
\textbf{Implemented tasks:}
\begin{itemize}[leftmargin=1.2cm]
    \item Construction of a Two-Tower neural network with tabular and ESM branches,
    \item Weighted BCE training to address class imbalance,
    \item Full LOGO loop across all 9 genes,
    \item Bootstrap confidence interval estimation for AUPRC,
    \item Brier and ECE-based calibration diagnostics,
    \item Calibration curves (reliability diagrams) for selected held-out settings,
    \item Per-gene in silico mutagenesis heatmap generation for Two-Tower and RF.
\end{itemize}

\subsection{Month 3: Finalization and Translational Output Layer}
\textbf{Implemented tasks:}
\begin{itemize}[leftmargin=1.2cm]
    \item Retraining final Two-Tower and RF models on full labeled data,
    \item Training and persisting a logistic Platt calibrator,
    \item VUS extraction from full cohort by uncertain-significance filtering,
    \item ESM delta extraction on VUS cohort,
    \item Pathogenicity probability ranking and CSV export for both model families.
\end{itemize}

\section{Two-Tower Architecture and Training Design}
\subsection{Network Layout}
The Two-Tower model uses three blocks:
\begin{enumerate}[leftmargin=1.2cm]
    \item \textbf{Tabular tower}: Linear $\rightarrow$ ReLU $\rightarrow$ BatchNorm $\rightarrow$ Dropout(0.3)
    \item \textbf{ESM tower}: Linear $\rightarrow$ ReLU $\rightarrow$ BatchNorm $\rightarrow$ Dropout(0.3)
    \item \textbf{Fusion head}: Concatenation $\rightarrow$ Linear/ReLU/BatchNorm/Dropout $\rightarrow$ Sigmoid output
\end{enumerate}

\subsection{Optimization and Imbalance Handling}
\begin{itemize}[leftmargin=1.2cm]
    \item Loss: weighted binary cross-entropy (sample-wise weight by class frequency)
    \item Optimizer: AdamW with $\text{lr}=0.003$, weight decay $=10^{-4}$
    \item Epochs: 40 in baseline training loops
\end{itemize}

\subsection{Calibration Strategy}
Raw neural probabilities are clipped for numerical stability and passed to logistic calibration (Platt scaling). The calibrator is trained on model outputs and reused in inference for calibrated ranking.

\section{Evaluation Protocol and Metrics}
\subsection{Cross-Gene Validation Principle}
The LOGO protocol withholds one gene at a time, trains on the remaining genes, and evaluates on the held-out gene. This setup is stricter than random splitting and directly measures cross-gene transfer.

\subsection{Metrics Reported in Pipeline Outputs}
\begin{itemize}[leftmargin=1.2cm]
    \item AUPRC (primary ranking metric under imbalance)
    \item 95\% bootstrap CI for AUPRC
    \item Brier score (probability calibration quality)
    \item ECE (Expected Calibration Error)
\end{itemize}

\subsection{Per-Gene LOGO Results from \texttt{results/logo\_metrics.csv}}
\begin{longtable}{p{0.12\textwidth} p{0.20\textwidth} p{0.11\textwidth} p{0.23\textwidth} p{0.11\textwidth} p{0.11\textwidth}}
\toprule
\textbf{Gene} & \textbf{Model} & \textbf{AUPRC} & \textbf{95\% CI} & \textbf{Brier} & \textbf{ECE} \\
\midrule
MYH7 & Two-Tower Hybrid & 0.8636 & [0.8355, 0.8865] & 0.5957 & 0.6515 \\
MYH7 & Baseline RF & 0.9086 & [0.8805, 0.9306] & 0.1673 & 0.1827 \\
MYBPC3 & Two-Tower Hybrid & 0.6061 & [0.5443, 0.6655] & 0.3548 & 0.3444 \\
MYBPC3 & Baseline RF & 0.6412 & [0.5748, 0.7114] & 0.2731 & 0.2272 \\
TNNT2 & Two-Tower Hybrid & 0.9206 & [0.8688, 0.9535] & 0.1446 & 0.0440 \\
TNNT2 & Baseline RF & 0.9205 & [0.8523, 0.9624] & 0.1611 & 0.1583 \\
TNNI3 & Two-Tower Hybrid & 0.9420 & [0.8711, 0.9864] & 0.1082 & 0.0771 \\
TNNI3 & Baseline RF & 0.9323 & [0.8785, 0.9737] & 0.1424 & 0.1510 \\
TPM1 & Two-Tower Hybrid & 0.9034 & [0.8327, 0.9735] & 0.1239 & 0.0500 \\
TPM1 & Baseline RF & 0.9649 & [0.9286, 0.9906] & 0.1248 & 0.1793 \\
ACTC1 & Two-Tower Hybrid & 0.9400 & [0.8801, 0.9871] & 0.1217 & 0.0636 \\
ACTC1 & Baseline RF & 0.9664 & [0.9197, 0.9930] & 0.1364 & 0.1717 \\
MYL2 & Two-Tower Hybrid & 0.8742 & [0.7398, 0.9635] & 0.2207 & 0.1867 \\
MYL2 & Baseline RF & 0.8593 & [0.7544, 0.9622] & 0.1839 & 0.1158 \\
MYL3 & Two-Tower Hybrid & 0.8261 & [0.6544, 0.9328] & 0.2656 & 0.2712 \\
MYL3 & Baseline RF & 0.8460 & [0.7277, 0.9402] & 0.2263 & 0.1599 \\
TNNC1 & Two-Tower Hybrid & 0.8324 & [0.6671, 0.9622] & 0.2066 & 0.1670 \\
TNNC1 & Baseline RF & 0.9339 & [0.8481, 0.9947] & 0.1564 & 0.1270 \\
\bottomrule
\end{longtable}

\subsection{Interpretation of LOGO Table}
\begin{itemize}[leftmargin=1.2cm]
    \item Both models perform strongly on multiple genes, but generalization varies by target gene.
    \item Two-Tower and RF are close on TNNT2, indicating robust signal capture from both model families.
    \item MYBPC3 remains relatively more difficult, suggesting heterogeneity and a possible need for richer contextual priors.
    \item Calibration characteristics vary materially by gene and model, motivating gene-aware thresholding in future deployment policy.
\end{itemize}

\section{Calibration Analysis and Reliability}
\subsection{Why Calibration Matters}
In translational settings, rank ordering alone is insufficient. Probability estimates guide triage thresholds, review prioritization, and confidence communication. Miscalibration can lead to overconfident false positives or underprioritized true pathogenic candidates.

\subsection{Reported Calibration Indicators}
\begin{itemize}[leftmargin=1.2cm]
    \item \textbf{Brier score}: lower values indicate better mean squared probability accuracy.
    \item \textbf{ECE}: lower values indicate better agreement between predicted confidence and empirical outcomes.
\end{itemize}

\subsection{Visual Reliability Outputs}
The pipeline saves reliability plots for selected folds in \texttt{figures/}, for example:
\begin{itemize}[leftmargin=1.2cm]
    \item \texttt{figures/calibration\_plot\_m2\_Two-Tower\_Hybrid.png}
    \item \texttt{figures/calibration\_plot\_m2\_Baseline\_RF.png}
\end{itemize}

\section{In Silico Mutagenesis (ISM) Outputs}
\subsection{Purpose of ISM in This Pipeline}
The ISM stage generates position-by-residue risk landscapes by mutating residues across selected windows and scoring each synthetic substitution. This helps identify mutation-sensitive regions and mutation classes likely to elevate model-assigned pathogenicity.

\subsection{Computation Flow}
For each target gene and selected positions:
\begin{enumerate}[leftmargin=1.2cm]
    \item Construct mutant 11-mer sequence for each amino acid replacement.
    \item Recompute ESM delta embeddings for mutant-vs-wild-type pair.
    \item Update mutation-sensitive tabular fields (size, charge, Grantham, center-window attributes).
    \item Score with calibrated Two-Tower and RF heads.
    \item Save heatmaps under \texttt{figures/}.
\end{enumerate}

\subsection{Generated Heatmap Assets}
Representative outputs include:
\begin{itemize}[leftmargin=1.2cm]
    \item \texttt{ism\_landscape\_m2\_TwoTower\_MYH7.png}
    \item \texttt{ism\_landscape\_m2\_TwoTower\_TNNT2.png}
    \item \texttt{ism\_landscape\_m2\_RF\_MYBPC3.png}
    \item \texttt{ism\_landscape\_m2\_RF\_TNNC1.png}
\end{itemize}

\subsection{Interpretive Value}
These landscapes are intended as mechanistic prioritization aids rather than causal proof. They are most useful for highlighting sequence positions where broad substitution classes are consistently assigned elevated risk.

\section{External Benchmarking Against EVE}
\subsection{Benchmark Scope}
A dedicated benchmark module under \texttt{eve\_benchmark/} compares project model outputs with EVE-derived references at per-gene level.

\subsection{Repository-level Summary}
From project documentation and benchmark outputs:
\begin{itemize}[leftmargin=1.2cm]
    \item Mean Two-Tower AUPRC across nine genes: \textbf{0.8427}
    \item Mean EVE AUPRC across nine genes: \textbf{0.8079}
\end{itemize}

\subsection{Benchmark Artifacts}
\begin{itemize}[leftmargin=1.2cm]
    \item Script: \texttt{eve\_benchmark/scripts/evaluate\_eve.py}
    \item Results: \texttt{eve\_benchmark/results/eve\_comparison.csv}
    \item Figures: \texttt{eve\_benchmark/figures/roc\_comparison.png}, \texttt{eve\_benchmark/figures/pr\_comparison.png}
\end{itemize}

\section{VUS Re-stratification Pipeline and Outcomes}
\subsection{VUS Extraction Rule}
Month 3 filters full cohort records whose \texttt{clin\_sig} contains ``Uncertain'' (case-insensitive), generating the working VUS dataset for re-ranking.

\subsection{Inference Steps for VUS Cohort}
\begin{itemize}[leftmargin=1.2cm]
    \item Align VUS tabular features with training feature columns.
    \item Extract live ESM 650M delta embeddings for each VUS variant.
    \item Produce Two-Tower calibrated probabilities and RF probabilities.
    \item Sort descending by predicted pathogenicity and export ranked tables.
\end{itemize}

\subsection{Result Files and Downstream Usage}
\begin{itemize}[leftmargin=1.2cm]
    \item \texttt{results/VUS\_restratification\_two\_tower.csv}
    \item \texttt{results/VUS\_restratification\_rf.csv}
\end{itemize}

These outputs can be integrated into expert review queues where higher-probability variants are prioritized for functional studies, segregation analysis, or additional evidence consolidation.

\section{Reproducibility and Engineering Operations}
\subsection{Execution Order}
The intended phase order is:
\begin{enumerate}[leftmargin=1.2cm]
    \item \texttt{python scripts/execute\_month1.py}
    \item \texttt{python scripts/execute\_month2.py}
    \item \texttt{python scripts/execute\_month3.py}
\end{enumerate}

\subsection{Operational Dependencies}
\begin{itemize}[leftmargin=1.2cm]
    \item Python ML stack: NumPy, Pandas, Scikit-learn, PyTorch
    \item Transformer stack for ESM extraction: \texttt{transformers}
    \item Visualization stack: Matplotlib, Seaborn
\end{itemize}

\subsection{Compute Considerations}
\begin{itemize}[leftmargin=1.2cm]
    \item ESM-2 extraction and ISM are computationally intensive.
    \item GPU acceleration materially improves runtime for embedding-heavy stages.
    \item Persisted intermediate artifacts reduce rerun cost and improve reproducibility.
\end{itemize}

\section{Risk Register and Limitations}
\subsection{Current Limitations}
\begin{enumerate}[leftmargin=1.2cm]
    \item \textbf{Class imbalance}: pathogenic majority may still bias certain decision surfaces.
    \item \textbf{Gene heterogeneity}: held-out behavior varies by gene, indicating non-uniform transfer difficulty.
    \item \textbf{Calibration variability}: confidence quality differs across model/gene slices.
    \item \textbf{Data dependency}: outputs are constrained by quality and representativeness of curated input cohorts.
\end{enumerate}

\subsection{Operational Risks}
\begin{itemize}[leftmargin=1.2cm]
    \item Heavy external model dependencies can challenge reproducibility across constrained environments.
    \item Long-run embedding extraction pipelines are vulnerable to interruption without checkpointing.
    \item Clinical use requires strict governance and expert oversight beyond algorithmic scoring.
\end{itemize}

\section{Translational Readiness and Deployment Considerations}
\subsection{What Is Ready Today}
\begin{itemize}[leftmargin=1.2cm]
    \item Production-grade persisted models and calibrator,
    \item Structured output tables for triage workflows,
    \item Visual interpretability artifacts for expert discussion,
    \item Clear phase-wise scripts for reproducible reruns.
\end{itemize}

\subsection{What Is Needed Before Clinical Operationalization}
\begin{itemize}[leftmargin=1.2cm]
    \item Independent external validation cohorts,
    \item Prospective calibration monitoring,
    \item Formal threshold policy tuned to intended clinical decision pathway,
    \item Documentation alignment with institutional governance requirements.
\end{itemize}

\section{Future Work Roadmap}
\begin{enumerate}[leftmargin=1.2cm]
    \item Extend evaluation to independent and newly released curated cohorts.
    \item Add stronger uncertainty decomposition and disagreement analytics across model families.
    \item Build checkpointed/parallelized embedding extraction for operational scalability.
    \item Add packaging/containerization for one-command reproducibility.
    \item Develop reviewer-facing reporting templates linking rank outputs to supporting interpretive evidence.
    \item Explore ensemble policies that combine calibration quality and ranking stability constraints.
\end{enumerate}

\section{Conclusion}
This expanded report consolidates the HCM repository into a full technical narrative suitable for a detailed academic or engineering submission. The project demonstrates a practical and auditable approach to variant pathogenicity prioritization that balances leakage control, cross-gene validation rigor, representation learning, and translational output generation.

The current repository already includes the key ingredients of a robust pipeline: curated data assets, leakage-aware preprocessing, hybrid and baseline model families, calibration diagnostics, in silico mutagenesis landscapes, external benchmark comparisons, and ranked VUS outputs. Together, these establish a strong basis for iterative enhancement toward higher confidence clinical decision support.

\section*{Appendix A: Artifact Checklist}
\begin{longtable}{p{0.32\textwidth} p{0.28\textwidth} p{0.32\textwidth}}
\toprule
\textbf{Artifact} & \textbf{Type} & \textbf{Purpose} \\
\midrule
\texttt{scripts/execute\_month1.py} & Python script & Cleaning, leakage removal, embeddings, baseline LOGO RF \\
\texttt{scripts/execute\_month2.py} & Python script & Two-Tower training, LOGO metrics, calibration, ISM \\
\texttt{scripts/execute\_month3.py} & Python script & Final training and VUS re-stratification \\
\texttt{data/HCM\_labeled\_final.csv} & Dataset & Labeled training/evaluation cohort \\
\texttt{data/HCM\_all\_variants\_v2.csv} & Dataset & Full cohort for uncertain-significance filtering \\
\texttt{data/esm2\_delta\_embeddings.npy} & Numpy artifact & Cached sequence embedding deltas \\
\texttt{models/hcm\_final\_two\_tower\_model.pth} & Model weights & Final Two-Tower inference model \\
\texttt{models/hcm\_final\_rf\_model.joblib} & Model bundle & Final RF inference model \\
\texttt{models/hcm\_platt\_calibrator.pkl} & Calibration artifact & Probability calibration for Two-Tower \\
\texttt{results/logo\_metrics.csv} & Result table & Per-gene LOGO performance and calibration indicators \\
\texttt{results/VUS\_restratification\_two\_tower.csv} & Result table & Ranked VUS list from calibrated Two-Tower model \\
\texttt{results/VUS\_restratification\_rf.csv} & Result table & Ranked VUS list from RF model \\
\texttt{eve\_benchmark/results/eve\_comparison.csv} & Benchmark table & EVE versus project-model comparison summary \\
\texttt{figures/*.png} & Figures & Calibration and ISM visual outputs \\
\bottomrule
\end{longtable}

\section*{Appendix B: Representative Figure Inventory}
\begin{itemize}[leftmargin=1.2cm]
    \item Calibration: \texttt{calibration\_plot\_m2\_Two-Tower\_Hybrid.png}, \texttt{calibration\_plot\_m2\_Baseline\_RF.png}
    \item Two-Tower ISM: \texttt{ism\_landscape\_m2\_TwoTower\_MYH7.png}, \texttt{ism\_landscape\_m2\_TwoTower\_MYBPC3.png}, \texttt{ism\_landscape\_m2\_TwoTower\_TNNT2.png}
    \item RF ISM: \texttt{ism\_landscape\_m2\_RF\_MYH7.png}, \texttt{ism\_landscape\_m2\_RF\_MYBPC3.png}, \texttt{ism\_landscape\_m2\_RF\_TNNT2.png}
    \item EVE benchmark: \texttt{eve\_benchmark/figures/roc\_comparison.png}, \texttt{eve\_benchmark/figures/pr\_comparison.png}
\end{itemize}

\vfill
\textit{Expanded report generated from repository artifacts in \texttt{abhi8667/HCM}.}

\end{document}
